% Minimal ICLR/NeurIPS-like host document: shows the figure and the companion
% table in a 5.5in single column, with a camera-ready caption.
%   pdflatex -output-directory=figures figures/demo_paper.tex
\documentclass[10pt]{article}
\usepackage[T1]{fontenc}
\usepackage[utf8]{inputenc}
\usepackage[letterpaper,textwidth=5.5in,textheight=9in]{geometry}
\usepackage{mathptmx}
\usepackage{helvet}
\usepackage{courier}
\usepackage{amsmath}
\usepackage{xcolor}
\usepackage{booktabs}
\usepackage{tikz}
\usetikzlibrary{positioning,calc,arrows.meta}
\usepackage[font=small,labelfont=bf]{caption}

\begin{document}

\begin{figure}[t]
\centering
% ---------------------------------------------------------------------------
% Video-QA evaluation pipeline: every configuration knob, stage by stage.
% Drop-in TikZ figure for ICLR / NeurIPS (single column, 5.5in = 13.97cm).
% Native width 13.56cm; nothing is scaled, so all type keeps its stated size.
%
% Requires in the host preamble:
%   \usepackage{tikz}\usetikzlibrary{positioning,calc,arrows.meta}
%   \usepackage{xcolor,amsmath}
% Usage:
%   \begin{figure}[t]\centering% ---------------------------------------------------------------------------
% Video-QA evaluation pipeline: every configuration knob, stage by stage.
% Drop-in TikZ figure for ICLR / NeurIPS (single column, 5.5in = 13.97cm).
% Native width 13.56cm; nothing is scaled, so all type keeps its stated size.
%
% Requires in the host preamble:
%   \usepackage{tikz}\usetikzlibrary{positioning,calc,arrows.meta}
%   \usepackage{xcolor,amsmath}
% Usage:
%   \begin{figure}[t]\centering% ---------------------------------------------------------------------------
% Video-QA evaluation pipeline: every configuration knob, stage by stage.
% Drop-in TikZ figure for ICLR / NeurIPS (single column, 5.5in = 13.97cm).
% Native width 13.56cm; nothing is scaled, so all type keeps its stated size.
%
% Requires in the host preamble:
%   \usepackage{tikz}\usetikzlibrary{positioning,calc,arrows.meta}
%   \usepackage{xcolor,amsmath}
% Usage:
%   \begin{figure}[t]\centering\input{figures/eval_pipeline}\caption{...}\end{figure}
% ---------------------------------------------------------------------------

% --- palette --------------------------------------------------------------
% Hues encode the sub-system; stage order and the accent rule carry the same
% information in grayscale, so the figure survives a b/w print.
\definecolor{cSrc}{RGB}{176,107,13}   % source & data
\definecolor{cVis}{RGB}{21,125,116}   % vision path
\definecolor{cTxt}{RGB}{94,75,192}    % text path
\definecolor{cGen}{RGB}{192,74,42}    % inference
\definecolor{cOut}{RGB}{94,122,16}    % scoring
\definecolor{cInk}{RGB}{28,27,31}
\definecolor{cInkTwo}{RGB}{95,92,102}

% --- helpers --------------------------------------------------------------
\providecommand{\kn}[1]{\texttt{#1}}
\providecommand{\warnmark}{{\setlength{\fboxsep}{.5pt}%
  \raisebox{.15ex}{\colorbox{cGen}{\textcolor{white}{%
  \fontsize{3.8}{3.8}\selectfont\bfseries\sffamily !}}}}}
\providecommand{\sep}{\,\textcolor{cInkTwo}{\textperiodcentered}\,\allowbreak}

\begingroup
\begin{tikzpicture}[
  font=\sffamily, x=1cm, y=1cm,
  stagenode/.style={
    rectangle, draw, line width=.35pt, rounded corners=.9pt,
    text width=3.12cm, minimum height=3.90cm, inner sep=0pt,
    anchor=north west, text=cInk},
  payload/.style={
    rectangle, draw=none, text width=3.12cm, align=left,
    inner xsep=2.6pt, inner ysep=2.2pt, anchor=north west,
    font=\ttfamily\fontsize{5.0}{6.0}\selectfont},
  flow/.style={-{Straight Barb[length=1.6pt,width=2.4pt]},
               line width=.5pt, draw=cInkTwo},
]

% ------------------------------------------------------------- geometry ---
\def\bw{3.12}     % box width
\def\dx{3.48}     % column pitch (box + 0.36 gap)  -> total 13.56cm
\def\bh{3.90}     % box height (uniform, so payload strips align)
\def\rowA{0}
\def\rowB{-4.95}
\def\midY{-4.61}  % routing lane for the row-wrap cable

% -------------------------------------------------------------- \stagebox --
% #1 id  #2 x  #3 y  #4 accent  #5 stage no.  #6 title  #7 source location
% #8 knob body  #9 payload leaving the stage
\newcommand{\stagebox}[9]{%
  \node[stagenode, fill=#4!6, draw=#4!50] (#1) at (#2,#3) {%
    \begin{minipage}[t][\bh cm][t]{2.80cm}
      \raggedright\setlength{\parindent}{0pt}%
      \vspace{3.4pt}
      {\fontsize{4.6}{4.6}\selectfont\bfseries\color{#4}STAGE~#5}\\[1.8pt]
      {\fontsize{6.8}{7.4}\selectfont\bfseries #6}\\[1.6pt]
      {\ttfamily\fontsize{4.8}{5.6}\selectfont\color{cInkTwo}#7}\\[3.0pt]
      {\fontsize{5.5}{6.7}\selectfont\color{cInk}#8}
    \end{minipage}};
  \fill[#4] ([xshift=.35pt,yshift=-.35pt]#1.north west)
    rectangle ([xshift=-.35pt,yshift=-1.5pt]#1.north east);
  \node[payload, fill=#4!16, text=#4!60!black] (#1out)
    at ([yshift=-1.6pt]#1.south west) {\textbf{$\rightarrow$}~#9};
}

% ============================================================== ROW 1 ======
\stagebox{s0}{0}{\rowA}{cSrc}{0}{Dataset row}{tasks/videomme/utils.py}{%
  \kn{-{}-tasks}\sep picks \kn{\_w\_subtitle}, \kn{\_long}, \kn{\_no\_visual}\\[2.2pt]
  \kn{data\_files}\sep 200-row local parquet vs.\ 2{,}700 upstream\\[2.2pt]
  \kn{-{}-limit}, \kn{-{}-offset}, \kn{test\_split}, \kn{HF\_HOME}%
}{doc + /path/<id>.mp4}

\stagebox{s1}{\dx}{\rowA}{cVis}{1}{Temporal sampling}{qwen\_vl\_utils:smart\_nframes}{%
  \kn{nframes}\,$=32$\sep fixed count, evenly spaced over the whole clip\\[2.2pt]
  \kn{fps}\,\warnmark\sep mutually exclusive: setting it drops \kn{nframes}\\[2.2pt]
  \kn{FRAME\_FACTOR}\,$=2$\,\warnmark\sep counts round down to even%
}{N x H0 x W0 x 3}

\stagebox{s2}{2*\dx}{\rowA}{cVis}{2}{Spatial resize}{qwen\_vl\_utils:smart\_resize}{%
  \kn{max\_pixels}\,$=301\,056$\,\warnmark\sep hard-clamped at $602\,112$ for video\\[2.2pt]
  \kn{min\_pixels}\,$=200\,704$\sep smaller frames are upscaled\\[2.2pt]
  sides snap to multiples of 28%
}{N x H x W (mod 28)}

\stagebox{s3}{3*\dx}{\rowA}{cVis}{3}{Serialization}{protocol.py:to\_openai\_msgs}{%
  \kn{pass\_video\_url}\,$=$\,\kn{False}\,\warnmark\sep N \kn{<|image\_pad|>} blocks, not one \kn{<|video\_pad|>}\\[2.2pt]
  not reachable from \kn{-{}-model vllm}\\[2.2pt]
  \kn{LMMS\_IMAGE\_\allowbreak ENCODE\_FORMAT}\ $=$\ \kn{PNG}%
}{N x base64 image\_url}

% ============================================================== ROW 2 ======
\stagebox{s4}{0}{\rowB}{cTxt}{4}{Prompt construction}{doc\_to\_text + kwargs}{%
  \kn{post\_prompt}\sep ``Answer with the option's letter \dots\ directly.''\\[2.2pt]
  subtitle preamble\,\warnmark\sep names subtitles even when none are passed\\[2.2pt]
  profiles keyed by the literal \kn{-{}-model} string%
}{messages[N imgs, text]}

\stagebox{s5}{\dx}{\rowB}{cGen}{5}{Context budget}{chat template + vLLM engine}{%
  \kn{max\_model\_len}\,$=32\,768$\sep too small truncates silently\\[2.2pt]
  \kn{limit\_mm\_per\_prompt}\,\warnmark\sep \kn{\{"video":1\}} caps a modality that never reaches the engine\\[2.2pt]
  \kn{gpu\_memory\_\allowbreak utilization}\ $=0.8$%
}{$\approx$6144 vision tok}

\stagebox{s6}{2*\dx}{\rowB}{cGen}{6}{Generation}{simple/vllm.py:sampling}{%
  \kn{max\_new\_tokens}\,\warnmark\sep $\max$(task~16, model~4096) is a \emph{floor}, not a cap\\[2.2pt]
  \kn{temperature}\,$=0$, \kn{top\_p}\,$=0.95$\\[2.2pt]
  \kn{num\_beams}, \kn{do\_sample} are dropped in silence%
}{"A" or essay $\supset$ "A"}

\stagebox{s7}{3*\dx}{\rowB}{cOut}{7}{Parse \& score}{process\_results + metrics}{%
  \kn{-{}-reasoning\_tags}\sep \kn{<think>} stripped before matching\\[2.2pt]
  extraction regex is per task\sep bare letter vs.\ \kn{mcq\_acc}\\[2.2pt]
  \kn{cluster\_key}\,$=$\,\kn{videoID} widens the SE, correctly%
}{acc $=0.6150$ $\pm$ SE}

% ============================================================== CABLES =====
\draw[flow] (s0out.east) -- (s1.west |- s0out.east);
\draw[flow] (s1out.east) -- (s2.west |- s1out.east);
\draw[flow] (s2out.east) -- (s3.west |- s2out.east);
\draw[flow] (s4out.east) -- (s5.west |- s4out.east);
\draw[flow] (s5out.east) -- (s6.west |- s5out.east);
\draw[flow] (s6out.east) -- (s7.west |- s6out.east);
\draw[flow, rounded corners=2.4pt]
  (s3out.south) -- (s3out.south |- 0,\midY)
                -- (s4.north  |- 0,\midY) -- (s4.north);

% ============================================================== FOOTER =====
\node[anchor=north west, inner sep=0pt] at (0,-9.62) {%
  \begin{minipage}{13.56cm}
  \raggedright\setlength{\parindent}{0pt}%
  \textcolor{cInkTwo}{\rule{13.56cm}{.35pt}}\\[4.0pt]
  {\fontsize{6.0}{7.6}\selectfont\sffamily
   \makebox[1.5em][l]{\warnmark}%
   \textbf{Knobs that do not do what their name suggests.}~%
   \kn{fps} and \kn{nframes} are mutually exclusive and never merged\sep
   \kn{max\_frame\_num} is read only on the \kn{fps} branch\sep
   odd frame counts round down (ask 7, get 6)\sep
   \kn{max\_pixels} is silently clamped to $602\,112$ for video frames\sep
   \kn{pass\_video\_url} cannot be reached from \kn{-{}-model vllm}\sep
   \kn{max\_new\_tokens} is a floor, so \kn{-{}-gen\_kwargs} alone will not lower it.\\[4.5pt]
   \makebox[1.5em][l]{\textcolor{cVis}{\bfseries$\Sigma$}}%
   \textbf{Stages 1--2 fix the context bill paid at stage 5.}~%
   $\text{vision\_tokens} \approx (\text{frames}/2)\times(\text{max\_pixels}/784)$,
   where the $/2$ is the temporal patch merge and $784 = 28\!\times\!28$ pixels per token
   after the $2\!\times\!2$ spatial merge; 32 frames at \kn{max\_pixels}\,$=301\,056$
   is ${\approx}6{,}144$ vision tokens for a single multiple-choice question.}
  \end{minipage}};

\end{tikzpicture}
\endgroup
\caption{...}\end{figure}
% ---------------------------------------------------------------------------

% --- palette --------------------------------------------------------------
% Hues encode the sub-system; stage order and the accent rule carry the same
% information in grayscale, so the figure survives a b/w print.
\definecolor{cSrc}{RGB}{176,107,13}   % source & data
\definecolor{cVis}{RGB}{21,125,116}   % vision path
\definecolor{cTxt}{RGB}{94,75,192}    % text path
\definecolor{cGen}{RGB}{192,74,42}    % inference
\definecolor{cOut}{RGB}{94,122,16}    % scoring
\definecolor{cInk}{RGB}{28,27,31}
\definecolor{cInkTwo}{RGB}{95,92,102}

% --- helpers --------------------------------------------------------------
\providecommand{\kn}[1]{\texttt{#1}}
\providecommand{\warnmark}{{\setlength{\fboxsep}{.5pt}%
  \raisebox{.15ex}{\colorbox{cGen}{\textcolor{white}{%
  \fontsize{3.8}{3.8}\selectfont\bfseries\sffamily !}}}}}
\providecommand{\sep}{\,\textcolor{cInkTwo}{\textperiodcentered}\,\allowbreak}

\begingroup
\begin{tikzpicture}[
  font=\sffamily, x=1cm, y=1cm,
  stagenode/.style={
    rectangle, draw, line width=.35pt, rounded corners=.9pt,
    text width=3.12cm, minimum height=3.90cm, inner sep=0pt,
    anchor=north west, text=cInk},
  payload/.style={
    rectangle, draw=none, text width=3.12cm, align=left,
    inner xsep=2.6pt, inner ysep=2.2pt, anchor=north west,
    font=\ttfamily\fontsize{5.0}{6.0}\selectfont},
  flow/.style={-{Straight Barb[length=1.6pt,width=2.4pt]},
               line width=.5pt, draw=cInkTwo},
]

% ------------------------------------------------------------- geometry ---
\def\bw{3.12}     % box width
\def\dx{3.48}     % column pitch (box + 0.36 gap)  -> total 13.56cm
\def\bh{3.90}     % box height (uniform, so payload strips align)
\def\rowA{0}
\def\rowB{-4.95}
\def\midY{-4.61}  % routing lane for the row-wrap cable

% -------------------------------------------------------------- \stagebox --
% #1 id  #2 x  #3 y  #4 accent  #5 stage no.  #6 title  #7 source location
% #8 knob body  #9 payload leaving the stage
\newcommand{\stagebox}[9]{%
  \node[stagenode, fill=#4!6, draw=#4!50] (#1) at (#2,#3) {%
    \begin{minipage}[t][\bh cm][t]{2.80cm}
      \raggedright\setlength{\parindent}{0pt}%
      \vspace{3.4pt}
      {\fontsize{4.6}{4.6}\selectfont\bfseries\color{#4}STAGE~#5}\\[1.8pt]
      {\fontsize{6.8}{7.4}\selectfont\bfseries #6}\\[1.6pt]
      {\ttfamily\fontsize{4.8}{5.6}\selectfont\color{cInkTwo}#7}\\[3.0pt]
      {\fontsize{5.5}{6.7}\selectfont\color{cInk}#8}
    \end{minipage}};
  \fill[#4] ([xshift=.35pt,yshift=-.35pt]#1.north west)
    rectangle ([xshift=-.35pt,yshift=-1.5pt]#1.north east);
  \node[payload, fill=#4!16, text=#4!60!black] (#1out)
    at ([yshift=-1.6pt]#1.south west) {\textbf{$\rightarrow$}~#9};
}

% ============================================================== ROW 1 ======
\stagebox{s0}{0}{\rowA}{cSrc}{0}{Dataset row}{tasks/videomme/utils.py}{%
  \kn{-{}-tasks}\sep picks \kn{\_w\_subtitle}, \kn{\_long}, \kn{\_no\_visual}\\[2.2pt]
  \kn{data\_files}\sep 200-row local parquet vs.\ 2{,}700 upstream\\[2.2pt]
  \kn{-{}-limit}, \kn{-{}-offset}, \kn{test\_split}, \kn{HF\_HOME}%
}{doc + /path/<id>.mp4}

\stagebox{s1}{\dx}{\rowA}{cVis}{1}{Temporal sampling}{qwen\_vl\_utils:smart\_nframes}{%
  \kn{nframes}\,$=32$\sep fixed count, evenly spaced over the whole clip\\[2.2pt]
  \kn{fps}\,\warnmark\sep mutually exclusive: setting it drops \kn{nframes}\\[2.2pt]
  \kn{FRAME\_FACTOR}\,$=2$\,\warnmark\sep counts round down to even%
}{N x H0 x W0 x 3}

\stagebox{s2}{2*\dx}{\rowA}{cVis}{2}{Spatial resize}{qwen\_vl\_utils:smart\_resize}{%
  \kn{max\_pixels}\,$=301\,056$\,\warnmark\sep hard-clamped at $602\,112$ for video\\[2.2pt]
  \kn{min\_pixels}\,$=200\,704$\sep smaller frames are upscaled\\[2.2pt]
  sides snap to multiples of 28%
}{N x H x W (mod 28)}

\stagebox{s3}{3*\dx}{\rowA}{cVis}{3}{Serialization}{protocol.py:to\_openai\_msgs}{%
  \kn{pass\_video\_url}\,$=$\,\kn{False}\,\warnmark\sep N \kn{<|image\_pad|>} blocks, not one \kn{<|video\_pad|>}\\[2.2pt]
  not reachable from \kn{-{}-model vllm}\\[2.2pt]
  \kn{LMMS\_IMAGE\_\allowbreak ENCODE\_FORMAT}\ $=$\ \kn{PNG}%
}{N x base64 image\_url}

% ============================================================== ROW 2 ======
\stagebox{s4}{0}{\rowB}{cTxt}{4}{Prompt construction}{doc\_to\_text + kwargs}{%
  \kn{post\_prompt}\sep ``Answer with the option's letter \dots\ directly.''\\[2.2pt]
  subtitle preamble\,\warnmark\sep names subtitles even when none are passed\\[2.2pt]
  profiles keyed by the literal \kn{-{}-model} string%
}{messages[N imgs, text]}

\stagebox{s5}{\dx}{\rowB}{cGen}{5}{Context budget}{chat template + vLLM engine}{%
  \kn{max\_model\_len}\,$=32\,768$\sep too small truncates silently\\[2.2pt]
  \kn{limit\_mm\_per\_prompt}\,\warnmark\sep \kn{\{"video":1\}} caps a modality that never reaches the engine\\[2.2pt]
  \kn{gpu\_memory\_\allowbreak utilization}\ $=0.8$%
}{$\approx$6144 vision tok}

\stagebox{s6}{2*\dx}{\rowB}{cGen}{6}{Generation}{simple/vllm.py:sampling}{%
  \kn{max\_new\_tokens}\,\warnmark\sep $\max$(task~16, model~4096) is a \emph{floor}, not a cap\\[2.2pt]
  \kn{temperature}\,$=0$, \kn{top\_p}\,$=0.95$\\[2.2pt]
  \kn{num\_beams}, \kn{do\_sample} are dropped in silence%
}{"A" or essay $\supset$ "A"}

\stagebox{s7}{3*\dx}{\rowB}{cOut}{7}{Parse \& score}{process\_results + metrics}{%
  \kn{-{}-reasoning\_tags}\sep \kn{<think>} stripped before matching\\[2.2pt]
  extraction regex is per task\sep bare letter vs.\ \kn{mcq\_acc}\\[2.2pt]
  \kn{cluster\_key}\,$=$\,\kn{videoID} widens the SE, correctly%
}{acc $=0.6150$ $\pm$ SE}

% ============================================================== CABLES =====
\draw[flow] (s0out.east) -- (s1.west |- s0out.east);
\draw[flow] (s1out.east) -- (s2.west |- s1out.east);
\draw[flow] (s2out.east) -- (s3.west |- s2out.east);
\draw[flow] (s4out.east) -- (s5.west |- s4out.east);
\draw[flow] (s5out.east) -- (s6.west |- s5out.east);
\draw[flow] (s6out.east) -- (s7.west |- s6out.east);
\draw[flow, rounded corners=2.4pt]
  (s3out.south) -- (s3out.south |- 0,\midY)
                -- (s4.north  |- 0,\midY) -- (s4.north);

% ============================================================== FOOTER =====
\node[anchor=north west, inner sep=0pt] at (0,-9.62) {%
  \begin{minipage}{13.56cm}
  \raggedright\setlength{\parindent}{0pt}%
  \textcolor{cInkTwo}{\rule{13.56cm}{.35pt}}\\[4.0pt]
  {\fontsize{6.0}{7.6}\selectfont\sffamily
   \makebox[1.5em][l]{\warnmark}%
   \textbf{Knobs that do not do what their name suggests.}~%
   \kn{fps} and \kn{nframes} are mutually exclusive and never merged\sep
   \kn{max\_frame\_num} is read only on the \kn{fps} branch\sep
   odd frame counts round down (ask 7, get 6)\sep
   \kn{max\_pixels} is silently clamped to $602\,112$ for video frames\sep
   \kn{pass\_video\_url} cannot be reached from \kn{-{}-model vllm}\sep
   \kn{max\_new\_tokens} is a floor, so \kn{-{}-gen\_kwargs} alone will not lower it.\\[4.5pt]
   \makebox[1.5em][l]{\textcolor{cVis}{\bfseries$\Sigma$}}%
   \textbf{Stages 1--2 fix the context bill paid at stage 5.}~%
   $\text{vision\_tokens} \approx (\text{frames}/2)\times(\text{max\_pixels}/784)$,
   where the $/2$ is the temporal patch merge and $784 = 28\!\times\!28$ pixels per token
   after the $2\!\times\!2$ spatial merge; 32 frames at \kn{max\_pixels}\,$=301\,056$
   is ${\approx}6{,}144$ vision tokens for a single multiple-choice question.}
  \end{minipage}};

\end{tikzpicture}
\endgroup
\caption{...}\end{figure}
% ---------------------------------------------------------------------------

% --- palette --------------------------------------------------------------
% Hues encode the sub-system; stage order and the accent rule carry the same
% information in grayscale, so the figure survives a b/w print.
\definecolor{cSrc}{RGB}{176,107,13}   % source & data
\definecolor{cVis}{RGB}{21,125,116}   % vision path
\definecolor{cTxt}{RGB}{94,75,192}    % text path
\definecolor{cGen}{RGB}{192,74,42}    % inference
\definecolor{cOut}{RGB}{94,122,16}    % scoring
\definecolor{cInk}{RGB}{28,27,31}
\definecolor{cInkTwo}{RGB}{95,92,102}

% --- helpers --------------------------------------------------------------
\providecommand{\kn}[1]{\texttt{#1}}
\providecommand{\warnmark}{{\setlength{\fboxsep}{.5pt}%
  \raisebox{.15ex}{\colorbox{cGen}{\textcolor{white}{%
  \fontsize{3.8}{3.8}\selectfont\bfseries\sffamily !}}}}}
\providecommand{\sep}{\,\textcolor{cInkTwo}{\textperiodcentered}\,\allowbreak}

\begingroup
\begin{tikzpicture}[
  font=\sffamily, x=1cm, y=1cm,
  stagenode/.style={
    rectangle, draw, line width=.35pt, rounded corners=.9pt,
    text width=3.12cm, minimum height=3.90cm, inner sep=0pt,
    anchor=north west, text=cInk},
  payload/.style={
    rectangle, draw=none, text width=3.12cm, align=left,
    inner xsep=2.6pt, inner ysep=2.2pt, anchor=north west,
    font=\ttfamily\fontsize{5.0}{6.0}\selectfont},
  flow/.style={-{Straight Barb[length=1.6pt,width=2.4pt]},
               line width=.5pt, draw=cInkTwo},
]

% ------------------------------------------------------------- geometry ---
\def\bw{3.12}     % box width
\def\dx{3.48}     % column pitch (box + 0.36 gap)  -> total 13.56cm
\def\bh{3.90}     % box height (uniform, so payload strips align)
\def\rowA{0}
\def\rowB{-4.95}
\def\midY{-4.61}  % routing lane for the row-wrap cable

% -------------------------------------------------------------- \stagebox --
% #1 id  #2 x  #3 y  #4 accent  #5 stage no.  #6 title  #7 source location
% #8 knob body  #9 payload leaving the stage
\newcommand{\stagebox}[9]{%
  \node[stagenode, fill=#4!6, draw=#4!50] (#1) at (#2,#3) {%
    \begin{minipage}[t][\bh cm][t]{2.80cm}
      \raggedright\setlength{\parindent}{0pt}%
      \vspace{3.4pt}
      {\fontsize{4.6}{4.6}\selectfont\bfseries\color{#4}STAGE~#5}\\[1.8pt]
      {\fontsize{6.8}{7.4}\selectfont\bfseries #6}\\[1.6pt]
      {\ttfamily\fontsize{4.8}{5.6}\selectfont\color{cInkTwo}#7}\\[3.0pt]
      {\fontsize{5.5}{6.7}\selectfont\color{cInk}#8}
    \end{minipage}};
  \fill[#4] ([xshift=.35pt,yshift=-.35pt]#1.north west)
    rectangle ([xshift=-.35pt,yshift=-1.5pt]#1.north east);
  \node[payload, fill=#4!16, text=#4!60!black] (#1out)
    at ([yshift=-1.6pt]#1.south west) {\textbf{$\rightarrow$}~#9};
}

% ============================================================== ROW 1 ======
\stagebox{s0}{0}{\rowA}{cSrc}{0}{Dataset row}{tasks/videomme/utils.py}{%
  \kn{-{}-tasks}\sep picks \kn{\_w\_subtitle}, \kn{\_long}, \kn{\_no\_visual}\\[2.2pt]
  \kn{data\_files}\sep 200-row local parquet vs.\ 2{,}700 upstream\\[2.2pt]
  \kn{-{}-limit}, \kn{-{}-offset}, \kn{test\_split}, \kn{HF\_HOME}%
}{doc + /path/<id>.mp4}

\stagebox{s1}{\dx}{\rowA}{cVis}{1}{Temporal sampling}{qwen\_vl\_utils:smart\_nframes}{%
  \kn{nframes}\,$=32$\sep fixed count, evenly spaced over the whole clip\\[2.2pt]
  \kn{fps}\,\warnmark\sep mutually exclusive: setting it drops \kn{nframes}\\[2.2pt]
  \kn{FRAME\_FACTOR}\,$=2$\,\warnmark\sep counts round down to even%
}{N x H0 x W0 x 3}

\stagebox{s2}{2*\dx}{\rowA}{cVis}{2}{Spatial resize}{qwen\_vl\_utils:smart\_resize}{%
  \kn{max\_pixels}\,$=301\,056$\,\warnmark\sep hard-clamped at $602\,112$ for video\\[2.2pt]
  \kn{min\_pixels}\,$=200\,704$\sep smaller frames are upscaled\\[2.2pt]
  sides snap to multiples of 28%
}{N x H x W (mod 28)}

\stagebox{s3}{3*\dx}{\rowA}{cVis}{3}{Serialization}{protocol.py:to\_openai\_msgs}{%
  \kn{pass\_video\_url}\,$=$\,\kn{False}\,\warnmark\sep N \kn{<|image\_pad|>} blocks, not one \kn{<|video\_pad|>}\\[2.2pt]
  not reachable from \kn{-{}-model vllm}\\[2.2pt]
  \kn{LMMS\_IMAGE\_\allowbreak ENCODE\_FORMAT}\ $=$\ \kn{PNG}%
}{N x base64 image\_url}

% ============================================================== ROW 2 ======
\stagebox{s4}{0}{\rowB}{cTxt}{4}{Prompt construction}{doc\_to\_text + kwargs}{%
  \kn{post\_prompt}\sep ``Answer with the option's letter \dots\ directly.''\\[2.2pt]
  subtitle preamble\,\warnmark\sep names subtitles even when none are passed\\[2.2pt]
  profiles keyed by the literal \kn{-{}-model} string%
}{messages[N imgs, text]}

\stagebox{s5}{\dx}{\rowB}{cGen}{5}{Context budget}{chat template + vLLM engine}{%
  \kn{max\_model\_len}\,$=32\,768$\sep too small truncates silently\\[2.2pt]
  \kn{limit\_mm\_per\_prompt}\,\warnmark\sep \kn{\{"video":1\}} caps a modality that never reaches the engine\\[2.2pt]
  \kn{gpu\_memory\_\allowbreak utilization}\ $=0.8$%
}{$\approx$6144 vision tok}

\stagebox{s6}{2*\dx}{\rowB}{cGen}{6}{Generation}{simple/vllm.py:sampling}{%
  \kn{max\_new\_tokens}\,\warnmark\sep $\max$(task~16, model~4096) is a \emph{floor}, not a cap\\[2.2pt]
  \kn{temperature}\,$=0$, \kn{top\_p}\,$=0.95$\\[2.2pt]
  \kn{num\_beams}, \kn{do\_sample} are dropped in silence%
}{"A" or essay $\supset$ "A"}

\stagebox{s7}{3*\dx}{\rowB}{cOut}{7}{Parse \& score}{process\_results + metrics}{%
  \kn{-{}-reasoning\_tags}\sep \kn{<think>} stripped before matching\\[2.2pt]
  extraction regex is per task\sep bare letter vs.\ \kn{mcq\_acc}\\[2.2pt]
  \kn{cluster\_key}\,$=$\,\kn{videoID} widens the SE, correctly%
}{acc $=0.6150$ $\pm$ SE}

% ============================================================== CABLES =====
\draw[flow] (s0out.east) -- (s1.west |- s0out.east);
\draw[flow] (s1out.east) -- (s2.west |- s1out.east);
\draw[flow] (s2out.east) -- (s3.west |- s2out.east);
\draw[flow] (s4out.east) -- (s5.west |- s4out.east);
\draw[flow] (s5out.east) -- (s6.west |- s5out.east);
\draw[flow] (s6out.east) -- (s7.west |- s6out.east);
\draw[flow, rounded corners=2.4pt]
  (s3out.south) -- (s3out.south |- 0,\midY)
                -- (s4.north  |- 0,\midY) -- (s4.north);

% ============================================================== FOOTER =====
\node[anchor=north west, inner sep=0pt] at (0,-9.62) {%
  \begin{minipage}{13.56cm}
  \raggedright\setlength{\parindent}{0pt}%
  \textcolor{cInkTwo}{\rule{13.56cm}{.35pt}}\\[4.0pt]
  {\fontsize{6.0}{7.6}\selectfont\sffamily
   \makebox[1.5em][l]{\warnmark}%
   \textbf{Knobs that do not do what their name suggests.}~%
   \kn{fps} and \kn{nframes} are mutually exclusive and never merged\sep
   \kn{max\_frame\_num} is read only on the \kn{fps} branch\sep
   odd frame counts round down (ask 7, get 6)\sep
   \kn{max\_pixels} is silently clamped to $602\,112$ for video frames\sep
   \kn{pass\_video\_url} cannot be reached from \kn{-{}-model vllm}\sep
   \kn{max\_new\_tokens} is a floor, so \kn{-{}-gen\_kwargs} alone will not lower it.\\[4.5pt]
   \makebox[1.5em][l]{\textcolor{cVis}{\bfseries$\Sigma$}}%
   \textbf{Stages 1--2 fix the context bill paid at stage 5.}~%
   $\text{vision\_tokens} \approx (\text{frames}/2)\times(\text{max\_pixels}/784)$,
   where the $/2$ is the temporal patch merge and $784 = 28\!\times\!28$ pixels per token
   after the $2\!\times\!2$ spatial merge; 32 frames at \kn{max\_pixels}\,$=301\,056$
   is ${\approx}6{,}144$ vision tokens for a single multiple-choice question.}
  \end{minipage}};

\end{tikzpicture}
\endgroup

\caption{\textbf{Everything between a video and a reported accuracy is
configuration.} A VideoMME row is an \texttt{mp4} path, a question, four options
and a gold letter; the eight stages above turn it into one number. Each panel
lists the settings that act at that stage, the module that reads them, and
(in the tinted strip) the representation leaving the stage. Markers~(\warnmark)
flag settings whose effect does not follow from their name: several are silently
ignored, silently clamped, or resolved as a floor rather than a cap. Frame count
and per-frame resolution are set entirely from \texttt{-{}-model\_args}, never
from the task YAML, and together they determine the context length paid at
stage~5. Values shown are those of our Qwen2.5-VL-7B-Instruct configuration
(vLLM 0.11.0, qwen-vl-utils 0.0.14, torchcodec reader).}
\label{fig:eval-pipeline}
\end{figure}

\begin{table}[t]
\centering
\caption{\textbf{Where an unexplained gap between two evaluation runs comes
from}, ordered by the size of the effect we observe. Stage numbers refer to
Figure~\ref{fig:eval-pipeline}.}
\label{tab:knobs}
% ---------------------------------------------------------------------------
% Companion appendix table for figures/eval_pipeline.tex.
% Full knob inventory, ordered by how much it moves the reported number.
% Requires: \usepackage{booktabs}
% Usage:  \begin{table}[t]\centering% ---------------------------------------------------------------------------
% Companion appendix table for figures/eval_pipeline.tex.
% Full knob inventory, ordered by how much it moves the reported number.
% Requires: \usepackage{booktabs}
% Usage:  \begin{table}[t]\centering% ---------------------------------------------------------------------------
% Companion appendix table for figures/eval_pipeline.tex.
% Full knob inventory, ordered by how much it moves the reported number.
% Requires: \usepackage{booktabs}
% Usage:  \begin{table}[t]\centering\input{figures/knob_table}\end{table}
% ---------------------------------------------------------------------------
\setlength{\tabcolsep}{4pt}
\renewcommand{\arraystretch}{1.12}
{\footnotesize
\begin{tabular}{@{}c l c p{0.505\linewidth}@{}}
\toprule
\# & Knob & Stage & Why it moves the number \\
\midrule
1  & \texttt{nframes} / \texttt{fps}                & 1 & Fewer frames carry less temporal evidence. Mutually exclusive: setting \texttt{fps} silently disables \texttt{nframes}. \\
2  & row set (\texttt{-{}-limit}, split, parquet)   & 0 & Different rows are not comparable at all; this checkout evaluates a 200-row subset, not the 2{,}700-row upstream split. \\
3  & \texttt{image\_pad} vs.\ \texttt{video\_pad}   & 3 & A genuinely different input representation, hence different logits. Not toggleable from \texttt{-{}-model vllm}. \\
4  & \texttt{max\_pixels} / \texttt{min\_pixels}    & 2 & Per-frame detail; silently clamped to $602\,112$ for video frames, so the requested value may not be the effective one. \\
5  & \texttt{max\_new\_tokens}                      & 6 & Resolved as $\max(\text{task},\text{model})$, a floor rather than a cap: 4096 instead of 16 lets the model ramble past the letter the parser expects. \\
6  & \texttt{FRAME\_FACTOR} rounding                & 1 & Frame counts round down to a multiple of two, so a run may use fewer frames than requested. \\
7  & prompt text / subtitles                        & 4 & Wording changes behaviour; the VideoMME preamble names subtitles even on the variant that passes none. \\
8  & \texttt{temperature}                           & 6 & Any value above zero makes the run non-reproducible. \\
9  & video reader backend                           & 1 & torchcodec, decord and torchvision land on marginally different frame indices. \\
10 & model snapshot / \texttt{chat\_template}       & 5 & Different weights or template change the tokenisation of every vision block. \\
\bottomrule
\end{tabular}}
\end{table}
% ---------------------------------------------------------------------------
\setlength{\tabcolsep}{4pt}
\renewcommand{\arraystretch}{1.12}
{\footnotesize
\begin{tabular}{@{}c l c p{0.505\linewidth}@{}}
\toprule
\# & Knob & Stage & Why it moves the number \\
\midrule
1  & \texttt{nframes} / \texttt{fps}                & 1 & Fewer frames carry less temporal evidence. Mutually exclusive: setting \texttt{fps} silently disables \texttt{nframes}. \\
2  & row set (\texttt{-{}-limit}, split, parquet)   & 0 & Different rows are not comparable at all; this checkout evaluates a 200-row subset, not the 2{,}700-row upstream split. \\
3  & \texttt{image\_pad} vs.\ \texttt{video\_pad}   & 3 & A genuinely different input representation, hence different logits. Not toggleable from \texttt{-{}-model vllm}. \\
4  & \texttt{max\_pixels} / \texttt{min\_pixels}    & 2 & Per-frame detail; silently clamped to $602\,112$ for video frames, so the requested value may not be the effective one. \\
5  & \texttt{max\_new\_tokens}                      & 6 & Resolved as $\max(\text{task},\text{model})$, a floor rather than a cap: 4096 instead of 16 lets the model ramble past the letter the parser expects. \\
6  & \texttt{FRAME\_FACTOR} rounding                & 1 & Frame counts round down to a multiple of two, so a run may use fewer frames than requested. \\
7  & prompt text / subtitles                        & 4 & Wording changes behaviour; the VideoMME preamble names subtitles even on the variant that passes none. \\
8  & \texttt{temperature}                           & 6 & Any value above zero makes the run non-reproducible. \\
9  & video reader backend                           & 1 & torchcodec, decord and torchvision land on marginally different frame indices. \\
10 & model snapshot / \texttt{chat\_template}       & 5 & Different weights or template change the tokenisation of every vision block. \\
\bottomrule
\end{tabular}}
\end{table}
% ---------------------------------------------------------------------------
\setlength{\tabcolsep}{4pt}
\renewcommand{\arraystretch}{1.12}
{\footnotesize
\begin{tabular}{@{}c l c p{0.505\linewidth}@{}}
\toprule
\# & Knob & Stage & Why it moves the number \\
\midrule
1  & \texttt{nframes} / \texttt{fps}                & 1 & Fewer frames carry less temporal evidence. Mutually exclusive: setting \texttt{fps} silently disables \texttt{nframes}. \\
2  & row set (\texttt{-{}-limit}, split, parquet)   & 0 & Different rows are not comparable at all; this checkout evaluates a 200-row subset, not the 2{,}700-row upstream split. \\
3  & \texttt{image\_pad} vs.\ \texttt{video\_pad}   & 3 & A genuinely different input representation, hence different logits. Not toggleable from \texttt{-{}-model vllm}. \\
4  & \texttt{max\_pixels} / \texttt{min\_pixels}    & 2 & Per-frame detail; silently clamped to $602\,112$ for video frames, so the requested value may not be the effective one. \\
5  & \texttt{max\_new\_tokens}                      & 6 & Resolved as $\max(\text{task},\text{model})$, a floor rather than a cap: 4096 instead of 16 lets the model ramble past the letter the parser expects. \\
6  & \texttt{FRAME\_FACTOR} rounding                & 1 & Frame counts round down to a multiple of two, so a run may use fewer frames than requested. \\
7  & prompt text / subtitles                        & 4 & Wording changes behaviour; the VideoMME preamble names subtitles even on the variant that passes none. \\
8  & \texttt{temperature}                           & 6 & Any value above zero makes the run non-reproducible. \\
9  & video reader backend                           & 1 & torchcodec, decord and torchvision land on marginally different frame indices. \\
10 & model snapshot / \texttt{chat\_template}       & 5 & Different weights or template change the tokenisation of every vision block. \\
\bottomrule
\end{tabular}}

\end{table}

\noindent Reported video-QA accuracies are rarely accompanied by the
configuration that produced them. Figure~\ref{fig:eval-pipeline} enumerates the
settings that a single VideoMME evaluation actually depends on, and
Table~\ref{tab:knobs} orders them by how far they move the reported number.

\end{document}
